\documentclass[11pt]{article}
\usepackage[margin=1in]{geometry}
\usepackage{amsmath,amssymb,amsthm,mathtools}
\usepackage{graphicx}
\usepackage{booktabs}
\usepackage[hidelinks]{hyperref}
\usepackage{microtype}
\setlength{\emergencystretch}{2em}

\newtheorem{theorem}{Theorem}
\newtheorem{corollary}[theorem]{Corollary}
\newtheorem{lemma}{Lemma}
\newcommand{\HH}{\mathsf H}
\newcommand{\PP}{\mathcal P_2}
\newcommand{\Wtwo}{W_2}
\newcommand{\Lip}{\operatorname{Lip}}
\newcommand{\Law}{\operatorname{Law}}
\newcommand{\E}{\mathbb E}

\title{Mixed-Sign Stochastic Neural Fields:\\
Full Exponential Ergodicity Without Kernel Definiteness}
\author{Automated mathematics run \texttt{fa\_banach\_001}, lane 1}
\date{26 August 2026}

\begin{document}
\maketitle

\begin{center}
\fbox{\parbox{0.91\textwidth}{
\textbf{Status: full resolution of the stated mixed-kernel exponential-
ergodicity problem, likely valid.}
In the source's quantitative ergodicity regime, its invariant-measure,
exponential-ergodicity, and exponential-mixing conclusions hold for arbitrary
mixed-sign kernels.  The sign and compact-subspace assumptions are
unnecessary.}}
\end{center}

\begin{abstract}
Otsetova and T\"olle identify the extension of their main exponential-
ergodicity result to kernels mixing excitation and inhibition as an open
problem.  We resolve that question in the same quantitative regime.  If
\[
 \kappa:=2\alpha-2\lVert K\rVert\Lip(f)-\Lip(B)^2>0,
\]
then the Markov semigroup contracts the quadratic Wasserstein metric by
\(e^{-\kappa t/2}\), has a unique invariant law with finite second moment, and
converges exponentially to it from every finite-second-moment initial law.
The proof uses synchronous coupling and a uniform It\^o moment estimate, so it
does not use any sign or symmetry of the kernel, nor the source's compact
nonlocal Hilbert subspace.  A direct corollary recovers the source's main
conclusions under the dimensionally consistent version of its Assumption 5,
with Assumptions 3 and 4 deleted.
\end{abstract}

\section{Source question}

The source is A.-M. Otsetova and J. M. T\"olle, \emph{Ergodicity for
stochastic neural field equations}, arXiv:2505.14012v1 (2025)\cite{source}.
On PDF page 5, after describing their sign-definite theory, they write:
\begin{quote}
``It would be interesting to obtain similar results, when the kernel is
neither non-negative or non-positive definite, relating to mixed inhibition
and excitation.''
\end{quote}

\begin{figure}[ht]
  \centering
  \includegraphics[width=\textwidth]{figures/open_problem_crop.png}
  \caption{The open direction in the source paper, PDF page 5.}
\end{figure}

Immediately before the open sentence, the source describes its main
contribution as existence of a unique invariant measure and exponential
ergodicity under a quantitative relation between decay, the kernel, the
activation, and the noise.  Thus the request concerns the exponential-
ergodicity regime of Theorem 1.2(iii); the following sentence identifies
metastability and Kramers' law as a separate open problem.

The source's existence argument restricts the dynamics to a compactly embedded
nonlocal Hilbert subspace built from a positive or negative part of the kernel.
That route is unavailable for a general mixed-sign operator.  The result below
replaces compactness by strict metric contraction.

\section{Setting and statement}

Let \(\HH=L^2(U,\rho)\), and consider the source equation
\begin{equation}\label{eq:sde}
 dX_t=(-\alpha X_t+K F(X_t))\,dt+B(X_t)\,dW_t,
 \qquad X_0=x\in\HH,
\end{equation}
where \(F(u)(\xi)=f(u(\xi))\).  We retain Assumptions 1 and 2 of
the source.  In particular, \(\HH\) is separable, \(K\in\mathcal L(\HH)\),
\(f:\mathbb R\to\mathbb R\) is Lipschitz, and
\(B:\HH\to L_2(V,\HH)\) is Lipschitz.  The source's Theorem 3.3 then gives a
unique global strong solution and its time-homogeneous Markov semigroup.

Set
\[
 L_f:=\Lip(f),\qquad L_B:=\Lip(B),\qquad
 \kappa:=2\alpha-2\lVert K\rVert L_f-L_B^2.
\]
No self-adjointness or quadratic-form sign is imposed on \(K\).

\begin{theorem}[Mixed-kernel contraction theorem]\label{thm:main}
Assume \(\kappa>0\).  Then:
\begin{enumerate}
\item for all \(\nu,\eta\in\PP(\HH)\) and \(t\ge0\),
\begin{equation}\label{eq:contraction}
 \Wtwo(P_t^*\nu,P_t^*\eta)
 \le e^{-\kappa t/2}\Wtwo(\nu,\eta);
\end{equation}
\item there is a unique invariant measure \(\mu\in\PP(\HH)\);
\item every \(\nu\in\PP(\HH)\) satisfies
\begin{equation}\label{eq:conv}
 \Wtwo(P_t^*\nu,\mu)
 \le e^{-\kappa t/2}\Wtwo(\nu,\mu);
\end{equation}
in particular, for every Lipschitz \(\phi:\HH\to\mathbb R\) for which the
integrals exist,
\[
 \left|P_t\phi(x)-\int_{\HH}\phi\,d\mu\right|
 \le \Lip(\phi)e^{-\kappa t/2}
 \left(\int_{\HH}\lVert x-y\rVert^2\,\mu(dy)\right)^{1/2}.
\]
\end{enumerate}
The conclusions hold for kernels with arbitrary mixtures of positive and
negative spectrum.
\end{theorem}

\begin{corollary}[Exact mixed-kernel extension of the source theorem]
\label{cor:source}
Retain source Assumptions 1 and 2, allow \(K\) to have arbitrary mixed
positive and negative spectrum, and write \(L_B\) for the Lipschitz constant
specified in source Assumption 2.  If the dimensionally consistent form of
the source's quantitative Assumption 5 holds,
\begin{equation}\label{eq:sourcecondition}
 2\sqrt2\,\lVert K\rVert L_f+L_B^2<2\alpha,
\end{equation}
then all conclusions of Theorem~\ref{thm:main} hold.  In particular, the
unique invariant law with second moment exists and the semigroup is
exponentially ergodic and exponentially mixing, without source Assumptions 3
and 4.
\end{corollary}

\begin{proof}
Condition \eqref{eq:sourcecondition} implies
\[
 2\alpha-2\lVert K\rVert L_f-L_B^2
 >2(\sqrt2-1)\lVert K\rVert L_f\ge0.
\]
Hence \(\kappa>0\), and Theorem~\ref{thm:main} applies.  The theorem's
\(W_2\) estimate implies both weak mixing and the bounded-Lipschitz
(Fortet--Mourier) estimate used in the source.
\end{proof}

\section{Proof intuition}

The sign assumption in the source is used to manufacture compactness, not to
control the distance between two solutions.  Drive two solutions by the same
Wiener process.  The linear decay contributes \(-2\alpha\) to the evolution of
their squared distance, while the neural interaction costs at most
\(2\lVert K\rVert L_f\) and the quadratic variation costs \(L_B^2\).
Thus \(\kappa>0\) gives a strict mean-square contraction.

Contraction alone gives uniqueness but not existence.  A second It\^o estimate
absorbs the affine terms \(f(0)\) and \(B(0)\), giving a second-moment bound
uniform in time.  The semigroup identity then makes \(P_t^*\delta_0\) a
Wasserstein-Cauchy family: evolve its earlier and later laws for their common
final time and contract the initial discrepancy.  Completeness of
\((\PP(\HH),\Wtwo)\) supplies the invariant limit.

\section{Proof}

\subsection{Synchronous contraction}

Let \(X_t^x\) and \(X_t^y\) solve \eqref{eq:sde} with initial conditions
\(x,y\), driven by the same cylindrical Wiener process, and write
\(D_t=X_t^x-X_t^y\).  Applying It\^o's formula to \(\lVert D_t\rVert^2\),
localizing the martingale if necessary, and taking expectations gives
\begin{align*}
 \frac{d}{dt}\E\lVert D_t\rVert^2
 &=-2\alpha\E\lVert D_t\rVert^2
 +2\E\langle D_t,K(F(X_t^x)-F(X_t^y))\rangle\\
 &\quad+\E\lVert B(X_t^x)-B(X_t^y)\rVert_{L_2(V,\HH)}^2.
\end{align*}
The Nemytskii map has Lipschitz constant \(L_f\).  Therefore
\[
 2\langle D_t,K(F(X_t^x)-F(X_t^y))\rangle
 \le 2\lVert K\rVert L_f\lVert D_t\rVert^2,
\]
and the noise term is bounded by \(L_B^2\lVert D_t\rVert^2\).  Hence
\[
 \frac{d}{dt}\E\lVert D_t\rVert^2
 \le-\kappa\E\lVert D_t\rVert^2,
\]
so Gr\"onwall's lemma yields
\begin{equation}\label{eq:pointcontract}
 \E\lVert X_t^x-X_t^y\rVert^2
 \le e^{-\kappa t}\lVert x-y\rVert^2.
\end{equation}

For any coupling of initial laws \(\nu,\eta\), use the coupled initial
conditions and the same Wiener process.  Taking the infimum over initial
couplings in \eqref{eq:pointcontract} proves \eqref{eq:contraction}.

\subsection{A uniform second moment}

All three alternatives in source Assumption 1 have
\(m_\rho:=\int_U\rho<\infty\).  Put
\[
 c_f:=|f(0)|\sqrt{m_\rho},\qquad b_0:=\lVert B(0)\rVert_{L_2(V,\HH)}.
\]
Then
\[
 \lVert F(u)\rVert\le L_f\lVert u\rVert+c_f,
 \qquad \lVert B(u)\rVert_{L_2}\le L_B\lVert u\rVert+b_0.
\]
Choose \(\varepsilon>0\) so that
\[
 q:=2\alpha-2\lVert K\rVert L_f-(1+\varepsilon)L_B^2>0;
\]
such a choice exists because \(\kappa>0\).  It\^o's formula and
\((a+b)^2\le(1+\varepsilon)a^2+(1+\varepsilon^{-1})b^2\) give, for
\(M_x(t):=\E\lVert X_t^x\rVert^2\),
\begin{align*}
 M_x'(t)
 &\le-qM_x(t)+2\lVert K\rVert c_f\sqrt{M_x(t)}
 +(1+\varepsilon^{-1})b_0^2\\
 &\le-\frac q2M_x(t)
 +\frac{2\lVert K\rVert^2c_f^2}{q}
 +(1+\varepsilon^{-1})b_0^2.
\end{align*}
Define
\[
 \gamma:=\frac q2,
 \qquad C_0:=\frac{2\lVert K\rVert^2c_f^2}{q}
 +(1+\varepsilon^{-1})b_0^2.
\]
Another application of Gr\"onwall gives
\begin{equation}\label{eq:moment}
 M_x(t)\le e^{-\gamma t}\lVert x\rVert^2+\frac{C_0}{\gamma}
 \qquad(t\ge0).
\end{equation}
In particular, the laws \(\nu_t:=P_t^*\delta_0\) have uniformly bounded
second moments.

\subsection{Existence, invariance, and uniqueness}

Let \(r\ge t\) and write \(r=t+s\).  By the semigroup property and
\eqref{eq:contraction},
\begin{align*}
 \Wtwo(\nu_r,\nu_t)
 &=\Wtwo(P_t^*\nu_s,P_t^*\delta_0)\\
 &\le e^{-\kappa t/2}\Wtwo(\nu_s,\delta_0)\\
 &\le e^{-\kappa t/2}
 \left(\sup_{a\ge0}\int_{\HH}\lVert z\rVert^2\,\nu_a(dz)\right)^{1/2}.
\end{align*}
The last supremum is finite by \eqref{eq:moment}.  Thus \((\nu_t)_{t\ge0}\)
is Cauchy as \(t\to\infty\) in \(\Wtwo\).  Since \(\HH\) is complete and
separable, \((\PP(\HH),\Wtwo)\) is complete\cite{Villani}; hence
\(\nu_t\to\mu\in\PP(\HH)\).

For fixed \(h\ge0\), contraction makes \(P_h^*\) continuous on
\(\PP(\HH)\), and
\[
 P_h^*\mu=\lim_{t\to\infty}P_h^*\nu_t
 =\lim_{t\to\infty}\nu_{t+h}=\mu.
\]
Thus \(\mu\) is invariant.  If \(\widetilde\mu\in\PP(\HH)\) is another
invariant law, then for every \(t>0\),
\[
 \Wtwo(\mu,\widetilde\mu)
 =\Wtwo(P_t^*\mu,P_t^*\widetilde\mu)
 \le e^{-\kappa t/2}\Wtwo(\mu,\widetilde\mu).
\]
It follows that \(\mu=\widetilde\mu\).  Applying
\eqref{eq:contraction} with \(\eta=\mu\) proves \eqref{eq:conv}.
Finally, Kantorovich--Rubinstein duality and \(W_1\le W_2\) give the stated
Lipschitz-observable estimate.  This completes the proof.

\section{Sharpness of the contraction coefficient}

The squared noise constant and the half-rate are not artifacts of the proof.
Let \(\HH=\mathbb R^2\), let
\[
 K=\begin{pmatrix}a&0\\0&-b\end{pmatrix},
 \qquad a\ge b>0,\qquad f(r)=r,
\]
and let \(B(u)q=\sigma q\,u\) for \(q\in\mathbb R\).  The kernel has genuinely
mixed spectrum, \(\lVert K\rVert=a\), and \(L_B=|\sigma|\).  If two solutions
start with a difference in the first eigendirection and are driven
synchronously, their difference \(D_t\) obeys
\[
 dD_t=(-\alpha+a)D_t\,dt+\sigma D_t\,dW_t.
\]
Consequently
\[
 \mathbb E\lVert D_t\rVert^2
 =e^{-(2\alpha-2a-\sigma^2)t}\lVert D_0\rVert^2.
\]
Thus \(\kappa=2\alpha-2\lVert K\rVert L_f-L_B^2\) is the exact
mean-square exponent already for a mixed-spectrum linear example.  For two
same-sign initial coordinates on this eigendirection, both laws are positive
scalar multiples of the same lognormal random variable; their monotone
coupling is optimal in one dimension.  Hence taking a square root gives the
exact \(W_2\) scale \(e^{-\kappa t/2}\).

\section{Verification and source comparison}

The proof uses only source Assumptions 1 and 2, the operator norm of \(K\),
and standard Hilbert-space It\^o calculus.  In particular, it never uses the
quadratic form \(\langle Ku,u\rangle\), so mixed positive and negative
spectrum causes no change.

Two constant conventions deserve attention.  The source defines its noise
constant by \(\lVert B(u)-B(v)\rVert_{L_2}\le C_B\lVert u-v\rVert\), but a
quadratic-variation estimate then naturally contains \(C_B^2\).  Also,
mean-square contraction at rate \(e^{-\kappa t}\) yields Wasserstein and
Lipschitz-observable contraction at rate \(e^{-\kappa t/2}\), after taking a
square root.  The statement here uses these squared/square-root conventions
consistently.  The separate verifier report records an adversarial step check.

\section{Resolution scope and novelty status}

This is a full answer to the source's stated mixed-kernel exponential-
ergodicity question in the quantitative regime in which the source formulates
its main ergodicity theorem.  Corollary~\ref{cor:source} recovers that theorem's
conclusions under its corrected smallness relation while deleting precisely
the two sign-dependent assumptions.  The result is stronger in two respects:
it assumes the sharper \(\kappa>0\) condition and gives \(W_2\) contraction.

No claim is made about invariant measures when \(\kappa\le0\), invariant laws
outside \(\PP(\HH)\), or the separately stated metastability and Kramers-law
problem.  These are broader parameter or moment questions, not qualifications
to the resolution of the source's strict-dissipation ergodicity question.

The run's four lightweight indexes and local parsed arXiv corpus were searched
using arXiv:2505.14012, the exact phrase about mixed inhibition and excitation,
the paper title, and close keywords concerning neural fields and Wasserstein
contraction.  Bounded web and arXiv searches used the same phrase and close
variants.  No separate paper explicitly resolving this source question was
found, including a check of the latest official arXiv source available on
26 August 2026.  The abstract synchronous-contraction mechanism is standard
in dissipative stochastic evolution equations, so the novelty claim is
deliberately limited to the exact sign-free resolution and its proof for this
source model.

\raggedright
\begin{thebibliography}{9}
\bibitem{source}
A.-M. Otsetova and J. M. T\"olle,
\emph{Ergodicity for stochastic neural field equations},
arXiv:2505.14012v1, 2025.

\bibitem{Villani}
C. Villani,
\emph{Optimal Transport: Old and New},
Grundlehren der mathematischen Wissenschaften 338, Springer, 2009.
\end{thebibliography}

\end{document}
